% ============================================================
% Part III — Desktop Section — M3 (The Deployment Record) — DRAFT v1
% Mode: Stage 3 guided drafting (per-paragraph protocol). M3.1 plan approved
%   2026-07-24; batch = M3.1 "The threshold" (3 ¶¶). Band R throughout;
%   hinge discipline: exactly ONE Band-G hinge sentence closes each
%   sub-module, naming the chain node; all argument deferred to M4.
% Conventions: as DESKTOP-M0-DRAFT-v2.tex / DESKTOP-M2-DRAFT-v1.tex (MLA;
%   respondent-code parentheticals for survey verbatims; chain nodes
%   A\textsubscript{n}; register + bans per memory).
% Data provenance: all rates/counts from t3-stats.md + 05-findings.md F1/F2
%   (all-Fable corpus, refreshed 2026-07-23). Verbatims verified char-exact
%   against data/open_text.csv 2026-07-24: W25-R042/Q9 (emoticon reproduced
%   verbatim — USER RULING 2026-07-24); W26-R059/Q9.
% Predicted-vs-unpredicted totals (¶2/¶3): predicted = lag 25 + nav 21 +
%   dark 10 + must-stand-close 2 = 58; unpredicted = crash 29 + install 26 +
%   hw 13 + A1.other 32 = 100. Excluded-at-threshold n = 61 (Q7-No ∧
%   video-only) = 25.3% of 241.
% Prediction-source framing: predictions registered in the coding design
%   (methodology-application/application-memo.md §4.2, fixed 2026-07-06,
%   before coding); Wendt is NOT named in section prose (not introduced by
%   the outline's M-structure; provenance lives in the analysis files).
% M3.1 — v2.1 APPROVED 2026-07-24 (LOCKED). M3.2 v1 DRAFTED 2026-07-24 —
%   IN REVIEW. Friction appendix DRAFT v1 written same day
%   (DESKTOP-APPENDIX-FRICTION-CATALOG-DRAFT-v1.tex; M3.1 ¶2 footnote
%   points to it). The M3.1 v1 rulings that produced v2 (style canon for
%   the rest of M3):
%   (a) REGISTER: too scientific — slow down, explain statistics in clearer,
%       more humanistic language; style sample presented for confirmation.
%   (b) NUMBERS: per-100-response rates unclear → switch to per-STUDENT
%       percentages (recomputed 2026-07-24 from 03-coded-corpus × respondents
%       q8 strata APP 123 / VID 108 / UNK 10): install 17% VID vs 3% APP;
%       hw 7% vs 3%; nonuse-talk 18% vs 0%; crash 15% APP vs 8% VID; lag 11%
%       vs 8%; nav 11% vs 6%. Raw counts stay. Define the two strata plainly
%       up front; "video route" → "YouTube route" (instrument's own wording:
%       "I utilized the YouTube videos only").
%   (c) ¶1: work through each obstacle in turn, one claim per sentence; the
%       user's model sentence for the Mac quote ADOPTED ("...excluded all
%       students who used Apple products, leading them to want it to be
%       '[m]ore inclusive for macbook users'").
%   (d) ¶2: grammar "the ones"→"those"; dev-team voice "we predicted five
%       elements..."; rename categories: VIDEO buffer latency (S3/AWS
%       streaming — delivery-system description noted for M2.1), video
%       fidelity/quality (not "dark-room"), close-proximity viewing; ALSO
%       predicted (AUTHOR-CONFIRMED 2026-07-24): startup crashes (UE5
%       performance-intensive), Mac-compatibility exclusion, missing YouTube
%       conveniences (playback speed omitted deliberately — pedagogical
%       decision). The 58-vs-100 predicted/unpredicted arithmetic is RETIRED.
%       Wendt DIRECT QUOTATION to include (user directive): breakdown-
%       anticipation, Design for Dasein pp. 33-34 (pdf 35-36), archon
%       exact-1.00+bbox verified: "anticipate the forms of breakdowns";
%       Bonsiepe qtd. in Wendt: breakdown "reveals the tissue of the
%       relations needed to fulfill the tasks".
%   (e) ¶3: de-stylize; thrust + final hinge sentence KEEP; tie in
%       world-disclosedness — the world FORECLOSED rather than disclosed for
%       technical reasons, not from any design aspect of the world itself.
%   (f) NEW APPENDIX ordered: every friction point, most→least reported
%       (scope proposal pending).
% ============================================================

\section{The Deployment Findings}
% [Title changed from "The Deployment Record" 2026-07-29, user-ordered at
%  assembly — "record" ban reaches the section title.]

% --- M3.1 The threshold (v2 REWRITE 2026-07-24, per approved style sample —
%     in review. Per-student percentages; strata defined in-prose; dev-team
%     prediction voice; Wendt quotation included; world-disclosedness close.
%     Data checks this session: 18 VID install-complainers, 15 of them Q7-No
%     ("problem was final" licensed); playback-speed request attested
%     W26-R021/Q9 + S25-R037/Q9; Wendt/Bonsiepe spans archon exact-1.00.)

The survey divides the 241 students into two groups, and the division itself
is the first finding. Each student reported whether they used the desktop
application or watched only the YouTube copies of the same footage: 123 used
the application, while 108---nearly half the class---chose the answer ``I
utilized the YouTube videos only.'' The YouTube group is itself two stories.
Sixty-one of its students reported that they could not use the application
at all; for them, YouTube was not a preference but the only door left open.
The other 47 reported that they could have used the application and turned
to YouTube anyway---a flight resulting from hassle rather than a failure to
access the platform.
The complaints make the difference between the groups concrete. Among the
108 students on the YouTube route, about one in six---17\%---described
problems installing the application or getting it to run, and for most of
them the problem was final: fifteen of those eighteen students are among the
sixty-one who never got in. Among the 123 application users, only 3\%
mention installation at all, and theirs are accounts of difficulty overcome,
since these students did, in the end, get the program running. The most
common obstacle was the operating system itself: the build ran only on
Windows, which excluded all students who used Apple products, leading them
to want it to be ``[m]ore inclusive for macbook users :('' (W25-R042/Q9).
The account gate stopped others before the hospital ever appeared; one
student's entire answer to the improvement question was ``Make sure the
sign-in is actually possible'' (W26-R059/Q9). A smaller number were stopped
by the machine itself, reporting hardware too weak for the application or no
Windows computer to put it on (7\% of the YouTube group, against 3\% of
application users).

The application users tell a different friction story, and theirs happens
inside the world. Their complaints are those the analysis was designed to
test. Given the development and testing of the Windows build, we predicted
five elements where students might encounter friction in using the virtual
space: buffering delay on the video streams, since every video is delivered
over the internet from cloud storage rather than shipped inside the
application; wayfinding, the work of locating a particular room in a
multi-story hospital; video fidelity, the sharpness of the footage once a
room's screen is playing it; the close proximity to the screen that
comfortable viewing requires; and avatar occlusion, one student's avatar
blocking another's view in a shared room. Predicting where a designed thing
will break is part of the work of design itself. The design theorist Thomas
Wendt, following Gui Bonsiepe, treats breakdowns as ``a necessary
accompaniment to design,'' each one revealing ``the tissue of the relations
needed to fulfill the tasks'' (qtd.\ in Wendt, \textit{Design for Dasein}
33--34); the designer's charge is to ``anticipate the forms of breakdowns''
rather than to imagine them away (\textit{Design for Dasein} 33--34). The
survey returned a verdict on each prediction. Buffering and latency drew
complaints from 11\% of application users (25 complaints in all);
wayfinding from another 11\% (21 complaints), often paired with
requests for maps or clearer guidance; video fidelity from 5\%, whose
complaints of blurry or hard-to-see footage name the worry in the students'
own vocabulary. The predicted proximity problem barely registers---two
complaints---and avatar occlusion never appears, for the simple reason that
most students entered the hospital alone: with no other avatars in the room,
there was no one to block the view. Beyond these five we had also
anticipated frictions of a blunter kind: crashes at startup, given how
demanding Unreal Engine 5 is of an ordinary student laptop; the exclusion of
Mac users; and the absence of conveniences students know from
YouTube---above all changing the playback speed, a feature we deliberately
withheld as a pedagogical decision. All three arrived. Crashes are the
largest single friction in the corpus (15\% of application users; 29
complaints); the Mac exclusion did its work at the threshold; and students
asked for exactly the feature we had withheld: ``I think it would be helpful
to have the option to change the playback speed of the videos''
(W26-R021/Q9).\footnote{A complete catalogue of every reported friction,
ordered from most to least reported, appears in the appendix.}
% [appendix number at assembly; catalogue draft ordered 2026-07-24]

Almost none of this concerns the world itself. Taken together, the two
friction stories name the machinery around the virtual hospital---the
installer, the operating system, the account gate, the stream, the frame
rate---and almost never its architecture, its rooms, or its footage as
designed. For the students it stopped, that machinery decided everything: a
world that cannot be reached discloses nothing, and what was lost at the
threshold was world-disclosedness itself. The virtual hospital foreclosed
rather than disclosed, and it foreclosed on technical grounds---not through
any design aspect of the world itself. In the chain's
terms the finding touches A\textsubscript{0} alone: availability is the
horizon's deployment face, and for a quarter of the class the world closed
at the installer, before anything in it could solicit so much as a look.

% --- M3.2 A world to see, nothing to do (v1 DRAFTED 2026-07-24 — in review.
%     ANCHOR F3; per-student stats computed 2026-07-24: D2-A3 166 apps /
%     147 students (61% of 241; APP 67% / VID 59%); D3-A3A4 37 apps /
%     33 students (15%/13%); G4-INTERACT 57 apps / 53 students (22%; APP 30%
%     vs VID 15%). Verbatims verified char-exact vs open_text.csv:
%     W25-R044/Q5+Q10, W25-R056/Q9, S25-R052/Q9, W26-R054/Q9 (student typo
%     "mor einteractive" — USER RULING 2026-07-24: quote verbatim incl.
%     typo, no [sic], per Mac-emoticon precedent), W25-R025/Q6, S25-R035/Q6,
%     W25-R058/Q6 (notes-outside-the-world anchor), W26-R027/Q9 ("occuring"
%     typo verbatim; inner double quotes → MLA single). Hinge = A3→A4.
%     orekton development EXPANDED per user ruling 2026-07-24 (ground in VLE
%     specifics + student responses). v2 of ¶2+¶3 written 2026-07-24:
%     same-student W25-R044 point made explicit; "One voice in 241" deleted;
%     "actually used the VLE" wording; action-talk defined in prose.)

Students praised what the world gave them to see. Mentions of the
content---the filmed operations, the physician interviews, the interest of
watching real surgery---are the largest positive channel in the whole
corpus: 166 coded mentions from 147 students, 61\% of the class. The praise
comes from both groups at nearly the same rate (67\% of application users,
59\% of the YouTube group), which follows from the deployment itself: both
routes carried the same footage, and what students valued was what the
footage showed. One student's summary can stand for the channel: ``the VR
simulation was super interesting and the wide-view videos of the procedures
being performed was very helpful'' (W25-R044/Q5).

Talk of doing is another matter. The coding distinguished students' talk
about the environment's content---what there was to watch---from talk about
acting in the environment: performing a task, making something, answering a
question, choosing a path. Talk of that second kind appears 37 times across
the 873 responses, from 33 students, and almost every instance is a request
or a complaint rather than a report: students describing not what they did
in the virtual hospital but what they wished it had given them to do. One
student put the wish in terms of routes: ``The VR clinical immersion
platform would be better if it was more interactable and had options on
which route you want to take rather than just being playback of videos''
(W25-R056/Q9). Another's entire answer to the improvement question was a
single request: ``Make assignments based on the VR clinical immersion''
(S25-R052/Q9). A third wrote: ``Maybe add something mor einteractive like
games/checklist/checkpoints etc'' (W26-R054/Q9). Requests of this kind came
from 53 students, 22\% of the class, and they came twice as often from
application users (30\%) as from students on the YouTube route (15\%): the
students asking for something to do are, above all, the students who
actually used the VLE. Among the 37 instances of action-talk, exactly one
reports an action the virtual hospital gave rather than an action it
lacked---and it comes from the same student, W25-R044, whose praise stood
for the content channel: ``The virtual space feels like you are in it,
while just a video doesn't. It makes you more involved and interactive with
the space'' (W25-R044/Q10).

What that missing act would have interrupted is long, unstructured
watching, and students describe that experience in detail. For a student
alone in the building---as most were---the world asks three things: find
the room, open its player, and watch. Nothing in the world asks anything
back. No room poses a question, no procedure ends in a prompt, and the
course's actual work of finding an unmet clinical need in the footage
happens outside the world entirely, in notes and documents the platform
does not include: ``I wish the VR platform would accommodate
other open windows. It wasn't easy to take notes on a Google doc while
listening to interviews or watching videos on the VR platform, so I
resorted to the youtube video afterwards'' (W25-R058/Q6). What remains
inside the world is watching, and the watching is long. One student found
``a lot of space of inactivity. Videos were unnecessarily long''
(W25-R025/Q6). Another asked for help with the seeking itself: ``I wish
there was a bit more guidance in watching the clinical procedures. Some of
them are many hours long, and I'm not sure what exactly to be looking for
during that time'' (S25-R035/Q6). A third met the same difficulty in the
course's own terms: ``at times it was hard to figure out a `problem' or
`unmet need' that was occuring. I think I was looking at the wrong things
in the beginning'' (W26-R027/Q9). These students are describing one
situation. Each arrived with attention to give and a task in hand---find a
need---and the footage gives every minute of the operation the same weight.
No moment is flagged as a difficulty; no prompt says \textit{watch how the
team works around this}; nothing distinguishes the places where an unmet
need shows itself---a workaround, a delay, an awkward instrument---from the
hours of routine that surround them. Seeking of the kind the course assigns requires a determinate
object---something in the world picked out as what the seeking is
\textit{for}---and the dissertation's term for that object is the
\textit{orekton}. The students quoted above report its absence in their own
words: interest solicited, seeking underway, nothing marked as the thing
sought. In the chain's terms the finding touches the passage from
A\textsubscript{3} to A\textsubscript{4}: the world reliably solicits
attention and interest, and reliably withholds the act that would complete
them.
% --- M3.3 The flight to video (v2 REVISED 2026-07-24 per user rulings —
%     in review. MAJOR THEORY RULING 2026-07-24: the flight is NOT a chain
%     stall — the chain COMPLETES by another route; object of desire (the
%     content) unchanged; doxa + emotion at A3 select the ROUTE. Hinge
%     recast to A3 (route-selection); developed analysis ¶4 added BY USER
%     ORDER (hinge-discipline relaxed here). v3 THEORY REFINEMENTS
%     2026-07-24: chain completes through SAME mechanisms always (never
%     "another route" for the chain itself — route varies only for object
%     delivery); phantasia expansion added per user — discursive-
%     deliberative orientational mode weighs both platform-phantasmata →
%     committal doxa (taking-as-true) in which emotion concretizes
%     (vocabulary verified against Part II v4 chain spine, line 231).
%     SEVEN-CAUSES ADDITION 2026-07-24 (user order): new ¶5 — voluntariness
%     cut maps the two anatomies (61 = compulsion, arche outside; 47 =
%     reasoning on habit's floor, arche within; "calculated retreat" echoes
%     Section 6's own example); habit kept strictly distinct from hexis per
%     standing ban; quotes = Rhet. I.10 1368b33-1369a6 + I.11 1370a6-8,
%     wording verified against Section 6 Live.md. "virtually natural"
%     ARCHON-VERIFIED 2026-07-24 exact-1.00+bbox vs Aristotle -
%     Rhetoric_(2014)_[Clean Copy].pdf pp.30-31 (cite I.11 1370a6-8 per
%     Section 6). ¶5 v2 rulings applied 2026-07-24: reasoning spelled out
%     (weighed steps vs return); "deliberate" not "calculated"; yes-or-no
%     STRUCTURE (also harmonized in ¶2); could-not/origin-within sentence
%     expanded with the giving-up student. ¶5 v3 2026-07-24: "comparative
%     reasoning enabled by phantasia" + DA III.11 434a5-10 footnote (verbatim
%     from Part I §1.3 / verbatim_passages, verified 2026-05-22 entry);
%     crossed-line sentence anchored in arche/compulsion/voluntary
%     vocabulary. Hinge extended with origin-within clause. Batch is 5 ¶¶.
%     M4 COORDINATION NOTE (user: recalibrate after phantasia-expansion
%     approval): the A3 route-selection analysis + seven-causes agency
%     individuation now live in M3.3 ¶¶4-5 — M4 must build on them, not
%     repeat them.
%     ANCHOR F4. Stats: B1-FALLBACK 48 apps /
%     41 students (VID 26% / APP 11%); anatomies 61 excluded vs 47 chosen
%     (Q7×Q8). Verbatims verified char-exact + strata checked 2026-07-24:
%     W25-R099/Q6 (Q7-No, excluded; "accomadate" typo verbatim),
%     W25-R017/Q9 (Q7-Yes∧VID — certified chosen flight), W25-R002/Q9
%     (Q7-Yes, APP), W25-R041/Q6 (Q7-No BUT narrates choice — soft-boundary
%     evidence, presented as such). Zeitvertreib grounded in prose per
%     M3.2 orekton precedent. YouTube-mirror deliberateness = M2.1 platform
%     fact (DEP), not re-cited.)

The 108 students who watched on YouTube divide into two groups. Sixty-one
reported that they could not run the application at all; for them the
flight was forced, and the videos were the accommodation the course had
built for exactly their case. Their talk is regret as often as complaint: ``I wish
that I could access the vr setting on macbook because it feels like a great
opportunity that I missed but the videos helped to accomadate for me''
(W25-R099/Q6). The other forty-seven students reported that they could have
used the application and watched on YouTube anyway. For them the world
stood open and they declined it. Talk of switching to the videos appears 48
times in the corpus, from 41 students, and it comes from 26\% of the
YouTube group against 11\% of application users---those last being students
who entered the world and then left it.

The students who declined the world explain themselves in the language of
effort. One, from the group that could have used the application, watched
the comparison tilt over time: ``the VR videos were cool at first, but the
youtube format ended up being simpler to use and watch'' (W25-R017/Q9). An
application user performed the same calculation from inside the world,
room by room: ``The effort going into finding a procedure to watch is way
more than just clicking on the youtube video, so it would be better if the
room had more than just the viewing feed and something to do to make that
effort worth it'' (W25-R002/Q9). A third student shows how soft the line
between the two anatomies really is. On the survey's yes-or-no item this
student answered that they could not use the platform, yet their account of
it is a decision, not a defeat: ``The software. I ended up just giving up
on it and just watching the youtube videos. It's not worth the effort to
even learn how to use it and download it all when I could just pull it up
on youtube'' (W25-R041/Q6). For some students, that is, ``could not'' did
not name a hard technical wall---the student was a Mac user and could not
run the program at all, or the hardware was too weak to run the build---but
the point at which they stopped trying: capable in principle of getting the application working, they
abandoned the download, the installation, or the troubleshooting when the
effort began to exceed what the platform seemed likely to repay. The line
between the sixty-one and the forty-seven is therefore softer than the
survey's yes-or-no structure suggests; an unknown share of the excluded students
excluded themselves. All three students quoted above are doing the same
calculation. Each step the virtual hospital requires---download, install,
sign in, launch, walk, find the room, open the player---is a step the
YouTube mirror does not, and what waited at the end of those steps was the
same footage. Friction sharpened the comparison where it appeared: mentions
of switching to the videos co-occur with install complaints eleven times,
with crashes five times, with latency four. A failed install or a crash
was, in case after case, the moment a student switched. But the moment is
not the mechanism. The mirror was available before any crash and after
every one---always one click away, always carrying the same footage, never
asking for a download, an account, or a walk through a building first.
Friction supplied the occasions; the mirror's standing cheapness supplied
the reason the occasion was always enough.

Watching on YouTube was never against the rules. The course created the
mirror channel deliberately, so that no student would be left without the
material, and for the sixty-one students whose machines could not run the
application---the Mac owners, the weak laptops, the failed installs---it
did exactly the job it was built for. But a channel built as an
accommodation for students who could not enter the world also stands as an
alternative for every student who could. Heidegger's name for the everyday
counter-measure to boredom is passing the time
(\textit{Zeitvertreib})---the reach for whatever occupation lies nearest
when the present one fails to hold us. The flight to video has exactly that
structure: not a verdict pronounced against the virtual hospital, but a
drift toward the nearest sufficient occupation---and here the drift was
institutional, because the course itself supplied the nearer channel and
stocked it with the same content. Whenever the virtual hospital offered a
student nothing beyond what the videos already offered---the same footage,
reached through more steps---the videos won, term after term, at a stable
43--46\% of each cohort.

The dissertation's frame makes the mechanics of that drift precise. The
flight is not a failure of the chain of desire, nor even an unusual run of
it: the chain completed through the same mechanisms by which it always
completes, and what varied was not the chain but the delivery of its
object. What these students desired---the object their coursework required
and their interest tracked---was the content itself: the operations, the
interviews, the material they needed to view. Nothing bound that object to
the virtual hospital; the same footage was available through both
distribution platforms. So the chain ran to its end. The footage solicited
interest. Between that interest and the act stood a comparison, and the
comparison is the work of \textit{phantasia}: in its
discursive-deliberative orientational mode, \textit{phantasia} holds a
\textit{phantasma} of each path before the mind---the download, the
install, the walk through the building on one side; the single click on the
other---and runs through routes not yet taken, which is what makes a
reasoned weighing of alternatives possible at all. That deliberation issued
in the committal \textit{doxa}, the taking-as-true: ``simpler to use and
watch'' (W25-R017/Q9); ``not worth the effort'' (W25-R041/Q6). In that
ratification the emotion concretized---the frustration of the threshold,
the impatience with the hassle---giving the judgment its weight. Then came
the act: the student watched, and the desire for the content was
actualized, on YouTube.

The \textit{Rhetorica}'s sorting of action by its causes names what kind of
act this was---and for whom. Aristotle divides all action by whether its
origin lies in the agent or outside: chance, nature, and compulsion move a
person from without; habit, reasoning, anger, and appetite are the agent's
own (\textit{Rhetorica} I.10, 1368b33--1369a6). The sixty-one students the
platform shut out fled by compulsion: the \textit{arch\=e} of their flight
was the machine's refusal, an external wall standing against their own
\textit{orexis}---the regret in ``a great opportunity that I missed''
(W25-R099/Q6) is the sound of a desire overridden from without. The
forty-seven who could have entered and did not fled from causes within
their own power, and their words name which. What their responses record is
comparative reasoning enabled by \textit{phantasia}: these students held
both platforms before the mind, measured them by a single standard---the
coursework's need for the content against the effort each platform
demanded---and pursued what appeared the greater good.\footnote{The
operation is the one Aristotle assigns to deliberative \textit{phantasia}:
``whether this or that shall be enacted is already a task requiring
calculation; and there must be a single standard to measure by, for that is
pursued which is greater. It follows that what acts in this way must be
able to make a unity out of several images'' (\textit{De Anima} III.11,
434a5--10).} Their retreat was a deliberate one. It rested, moreover,
on habit's floor: YouTube was for every student the
often-done thing, the channel whose use had long since become, in
Aristotle's phrase, ``virtually natural'' (\textit{Rhetorica} I.11,
1370a6--8)---so the familiar route asked nothing of them while the
unfamiliar one asked much. The taxonomy also says precisely what the
survey's yes-or-no structure blurred: the line between the two anatomies of
the flight is the line between an \textit{arch\=e} outside the student and
an \textit{arch\=e} within. Students like the one who ``ended up just
giving up on it'' answered that they could not use the platform, but their
own accounts locate the \textit{arch\=e} of the flight elsewhere. Had the
machine truly refused them---a build their hardware could not run, an
install that failed no matter how long they worked at it---the cause of
their flight would have been compulsion: an origin outside the agent,
external circumstance overriding the student's own desire to enter, the
involuntary cause under which the genuinely excluded students fall. What
these students describe instead is a stopping-point they chose. The
install had not finally failed; they had stopped installing, on their own
deliberation that the remaining effort was not worth its return. In the
\textit{Rhetorica}'s strict sense their flight was voluntary, the act their
own, whatever the survey answer records. The world, in the end, was not refused as
a world; it was weighed as a delivery platform for an object it did not
exclusively hold, and the effort it asked was too great for a destination
reachable in one click. In the chain's terms the finding touches A\textsubscript{3}: the
\textit{phantasmata} of both platforms taken up in the
discursive-deliberative mode and ratified by a \textit{doxa} in which the
emotion concretized---the chain completing as it always completes, deciding
not only whether the act would happen, but how, and from an origin that
was, for the forty-seven, entirely their own.
% --- M3.4 Presence happens on screens (v1 DRAFTED 2026-07-24 — in review.
%     ANCHOR F7 first half; hinge = A2 + world-disclosedness. 70.7% (87/123)
%     per standing correction, NOT the outline's ≈69%. Q8 item wording
%     quoted VERBATIM from raw Canvas header (W25 Student Analysis CSV,
%     item 3427111), same all terms per 00-audit. Pilot 70% verified
%     against p3-phenom-raw p. 388 ("70% of students noted the VR version
%     elicited a greater felt sense of presence"); cite form
%     ("Phenomenological Evaluation" 388) per M0 convention; pilot survey
%     respondent N=20 — prose avoids N to sidestep the 22-enrolled/20-
%     answered mismatch. F1-PRESENCE per-student 2026-07-24: 19 students,
%     APP 13 (11%) vs VID 6 (6%). G1-PROVISION 9 students; verbatims
%     W26-R014/Q9 + W25-R020/Q9 verified char-exact. W25-R044 same-student
%     continuity made explicit (per standing preference).
%     v2 REWRITE 2026-07-24 per user rulings: per-term percentages spelled
%     out; "deployed VLE"; scale = 20 vs 123 respondents stated; "through
%     traditional displays as much as 3D VR"; ¶2 expanded analytically —
%     world-disclosedness ("place dwelt in vs scene surveyed", Part II v4
%     L235 formula) + rhetorical incorporation (comes to full actuality at
%     committed-cognition→act; presence = "half its work") + phantasia
%     subtracts-the-surround (concept-lock framing); ¶3 opens direct;
%     pilot library-headset inversion FOOTNOTED (p. 389 verified, "not a
%     single student went there to check one out" verbatim);
%     TERMINOLOGY-LOCK COMPLIANCE (memory: rdr2-secondary-integration —
%     "phantasia-load"/"data-discrepancy"-in-analyses/"image-making" ALL
%     BANNED in analytical prose; earlier claim that v4 L205-207 licenses
%     them was WRONG for analyses — those are metaphysics-section coinages;
%     approved framing = medium-supply vs phantasia gap-filling /
%     subtract-the-surround): division-of-labor prose carries the concept
%     with phantasia only; hinge recast at the phantastic motion rendering
%     A2 (world taken-as dwelt in).
%     v3 EXPANSIONS 2026-07-24 (user rulings): same-display point added
%     (¶2 — difference is DESIGN not hardware; film can world-disclose;
%     phantasia active on both routes); incorporation-debt worked through
%     in own ¶ (disclosure = world's half of exchange; act ratifies;
%     begun-and-suspended; teleology = disclosure for action as means for
%     ends); NEW ¶ — matched-numbers insight (70% headsets vs 70.7%
%     monitors → immersive power from explorable-world inclusion, not
%     stereoscopic viewpoint) + Calleja: definition ARCHON-VERIFIED
%     exact-1.00+bbox (In-Game 169, pp.168-170 chunk) + spatial/kinesthetic
%     cornerstone pair (In-Game 170, digest ig-10 line-anchored); bare
%     "incorporation" for Calleja per standing rule. Batch is now 5 ¶¶.
%     Wendt considered for this insight, Calleja fit exactly (screen-game
%     provenance).)

Each term the survey put a direct question to the students who used the
application: ``Did you find the VR clinical immersion platform to elicit a
greater degree of immersion, presence, and/or embodiment (the feeling that
you are actually present in the operating rooms/clinical environment)?''
Of the 123 students who used the application and answered, 87 said
yes---70.7\%. The majority held in every term taken alone, not merely in
the pooled total: 64\% of Winter 2025's application users answered yes,
82\% of Spring 2025's, and 67\% of Winter 2026's. The figure lands with a
history behind it. In the pilot, we had asked whether the 3D VR versions
elicited a greater degree of presence than the other formats, and 70\% of
respondents said they did (``Phenomenological Evaluation''
388)---students wearing headsets, standing inside stereoscopic footage.
The deployed VLE reproduces the pilot's number almost exactly, with no
headset involved: where the pilot's question was answered by twenty
students in a single quarter, the deployment's immersion item was answered
by 123 application users across three terms, every one of them working
through an ordinary monitor by mouse and keyboard. Presence, as the
students report it, happened through traditional displays as much as it
did through 3D VR.

The yes-answers do not stand alone; the open responses corroborate them
in the students' own words. Nineteen students, unprompted by any
yes-or-no item, wrote descriptions of feeling present in the virtual
hospital---of being in the space rather than watching it---and thirteen
of the nineteen were application users (11\% of that group, against 6\%
of the YouTube group). The fullest of these descriptions comes from the
same student, W25-R044, who supplied the corpus's single account of given
action: ``The virtual space feels like you are in it, while just a video
doesn't'' (W25-R044/Q10). The sentence repays slow reading, because in a
student's plain words it draws the dissertation's own distinction---and
draws it under conditions that isolate what makes the difference. Both
routes reach this student through the same hardware: the YouTube videos
and the virtual hospital alike arrive on an ordinary computer display.
Whatever separates being-in it from watching it therefore cannot be the
screen; it is the design---how the same footage is framed and delivered
to the student who watches. Nor is traditional video barred from
disclosing a world: a film seen on that same monitor can take its viewer
into its world, and \textit{phantasia} is active on both routes. The
difference lies in what the design gives \textit{phantasia} to work
with. The YouTube route delivers the footage as one video among
videos---a page, a player, a list of what to watch next---content framed
as content, the surround reasserted at every edge of the frame. The VLE
delivers the same footage as the interior of a place: to reach it the
student has walked a building, found a room, and stood before its
screen, and the hospital continues around the footage while it plays.
What the student's clause registers is the difference between those two
relations---a world disclosed as a place dwelt in, against a scene
surveyed from without. That difference is what the dissertation names
world-disclosedness: the student taken into the world acted in, and the
world taken into the student's awareness. On a flat screen that
disclosure is \textit{phantasia}'s victory in either case. A monitor
gives the eyes a bounded rectangle inside a room that remains fully
present around it, and it is \textit{phantasia} that subtracts the
surround, supplies what the screen withholds, and encloses the student
in the world the screen only samples. The designs differ in how much
they help: the virtual hospital hands \textit{phantasia} a surrounding
place already built---an inside to be-in---where the video page hands it
a frame to be overcome. The coded responses register that difference in
help directly: of the nineteen students who described being-in, thirteen
came from the 123 application users and six from the 108 on the YouTube
route---students given the built world reported being-in at nearly twice
the rate of students given the framed video, though both watched the
same footage on the same kind of screen.

This disclosure is also the first face of what the dissertation calls
rhetorical incorporation: the actualization toward which a designed
world drives the whole chain of desire, and which comes to full
actuality only where committed cognition passes into a completed act.
It is worth working through what that means for these students. Feeling
present is not the completion of rhetorical incorporation; it is the world's half
of an exchange. The chain is understood backward from its end---from
the desired object whose pursuit completes it in an act---so a world
that has disclosed itself has, so far, only made acting possible: it
has given the student a place in which an act could happen. Rhetorical
incorporation is finished when the student's own act, done in and with
the world, ratifies the disclosure---when being-in becomes acting-in. A
world that discloses but does not task therefore leaves the exchange
half-done. The student stands in a place that asks nothing of them; the
disclosure that invites completion remains an invitation; rhetorical
incorporation is begun and then suspended at the very threshold where it
would become actual. That is the precise sense in which such a world still owes its
students an act: the debt is written into the chain's own teleology,
where disclosure is for action as means are for ends. Other students
put the being-in more modestly---``Being able to get somewhat first
hand observations'' (W25-R089/Q5); ``I found the immersive VR helpful''
(W26-R028/Q5)---but the relation they report is the same.

Nine students asked the course to provide VR headsets: ``This class
should provide students with VR headsets. Even if there were a technology
fee, it would be more beneficial for learning'' (W26-R014/Q9); ``I feel
like we should have some sort of cheap VR set because the VR experience
feels off when seeing it in the laptop screen''
(W25-R020/Q9).\footnote{The request inverts the pilot's most curious
result. In the pilot, we housed several Meta Quest 2 headsets in the
school's library, and ``not a single student went there to check one
out''; as one student phrased it, obtaining a headset ``was a hassle''
(``Phenomenological Evaluation'' 389). Three years on, students ask for
headsets to be handed to them. The desire for the richer medium is real
in both cohorts; what governs it is the same convenience gradient that
routes everything else in this deployment---intensity is wanted, but only
at zero fetching cost.} What these nine notice is the division of labor
just described. Every medium supplies some share of what the senses would
receive in the non-virtual scene, and \textit{phantasia} must furnish the
rest: a headset supplies much and asks little of \textit{phantasia}; a
flat screen supplies least and asks most. The student who writes that the
experience ``feels off\,\ldots\,in the laptop screen'' is reporting how
large \textit{phantasia}'s share becomes on the desktop route---feeling,
in the moment of viewing, how much of the world it falls to their own
\textit{phantasia} to hold up. What the 70.7\% then says is that the
share was held. Seven in ten application users, on the route where the
medium supplies least and \textit{phantasia} must supply most, reported
the feeling of actually being present in the operating room.

Set the pilot's number and the deployment's side by side, and they
sharpen into a finding about where immersion comes from. The pilot's
students wore stereoscopic headsets and stood inside the footage; the
deployment's students sat before flat monitors; both cohorts reported
presence at 70\%. If the stereoscopic viewpoint were what
carried immersion, the number should have fallen when the headsets
disappeared---and it did not fall. What the two cohorts shared was not
display hardware but a designed, explorable world: the virtual
hospital, entered and navigated, on whichever display. The likeliest
reading of the matched numbers is therefore that the immersive power
\textit{phantasia} conveys derives not from the stereoscopic viewpoint
but from the inclusion of the footage in an explorable world. Calleja's
account of incorporation, built from screen games---players at
monitors, no headset anywhere---predicts exactly this. He defines
incorporation as ``the absorption of a virtual environment into
consciousness, yielding a sense of habitation, which is supported by
the systemically upheld embodiment of the player in a single location,
as represented by the avatar'' (\textit{In-Game} 169): habitation grows
from emplacement in a navigable world, not from any property of the
display. Of the six dimensions of involvement his model blends into
incorporation, the cornerstone pair without which it cannot occur is
the \textit{spatial} and the \textit{kinesthetic}---the internalized
space and the learned movement through it (\textit{In-Game} 170)---and
that pair is
precisely what the VLE adds to the footage and YouTube does not. Behind
the mechanism stands a principle of Merleau-Ponty's: ``What counts for
the orientation of the spectacle is not my body, such as it in fact
exists, as a thing in objective space, but rather my body as a system
of possible actions, a virtual body whose phenomenal `place' is defined
by its task and by its situation. My body is wherever it has something
to do'' (\textit{Phenomenology of Perception} 260).\footnote{Wendt
reaches this same passage from the design side: for him it carries
direct consequences for the design of experience, which must orient
itself to what users are there to do---their goals, conscious and
non-conscious---within the multiple stable uses any technology admits
(\textit{Design for Dasein} 126). The convergence is the point: what
the phenomenologist states of the body, the design theorist prescribes
for the world built around it---a world is designed well when it gives
the virtual body something to do. A third tradition arrives at the same
structure independently. Chalmers, with no reference to Merleau-Ponty,
builds his defense of virtual realism on the same figure: ``If I pick
up a virtual coin in Second Life, I really use my virtual body to take
possession of a virtual coin. There is nothing fictional about this''
(``The Virtual and the Real'' 317); ``virtual actions are plausibly
real actions (albeit with a virtual body), so that when one performs
virtual actions, one really is doing something'' (339); and a user's
sense of ownership over a virtual body ``need not always be an
illusion'' (333). Phenomenology states the principle---the body is
where its tasks are; design prescribes it---build worlds that give the
virtual body something to do; analytic metaphysics certifies its
object---what the virtual body does is really done. For the survey's
70.7\% the triangulation matters: the students' reports of being
present in the operating room need not be read as a flattering
illusion, since they describe a virtual body genuinely located by task
and situation. It sharpens the deployment's deficit in the same stroke,
for a virtual body whose acts are real still needs acts to perform.}
% [Chalmers cites verified char-exact vs corpus PDF (pp. 317/333/339,
%  journal pagination); PDF marked 2016, published Disputatio 2017 —
%  settle year at Works Cited. ADD Chalmers TO WORKS CITED.]
% [CITE UPGRADED 2026-07-24: direct quotation from user's edition —
%  Merleau-Ponty, Phenomenology of Perception, Landes trans. (Routledge
%  2014 [1945]), p. 260, Part Two "Space" (Wertheimer passage); wording
%  verified via pdftotext (Landes: "wherever it has something to do",
%  NOT Smith's "wherever there is something to be done"). ADD TO WORKS
%  CITED. Wendt footnote = PARAPHRASE gloss (digest's "major
%  implications" phrase did NOT verify verbatim — no quotes used).
%  PENDING OPS: archon ingest of the new PhP PDF + snapshot rebuild
%  (vector-compact) queued for next ingestion session — OP RULE.
%  CORRECTION: Chalmers "Virtual and the Real" does NOT quote this
%  passage (earlier archon hit = partial-phrase match on his own
%  "virtual body" usage); his virtual-body realism converges but does
%  not cite M-P. CONTINUATION NOTED for possible later use (PhP
%  260-261): possible actions "sketch out...a possible habitat";
%  "the reflected room conjures up a subject capable of living in it.
%  This virtual body displaces the real body" — meets Calleja's
%  "sense of habitation" almost verbatim.]
The explorable hospital gives the student's virtual body somewhere to
be---a place constituted by situation and task, a room to find and a
player to open---on whatever display the world arrives through. In the
chain's terms the finding sits at the phantastic motion that renders
A\textsubscript{2}: from the thin \textit{aisth\=emata} a monitor
delivers, \textit{phantasia} still composes the charged
\textit{phantasma} of a surrounding world---a world taken-as dwelt
in---and it composes it from the explorable place the design supplies,
not from the apparatus on the head.
% --- M3.5 The WE: had, wanted, unactivated (v1 DRAFTED 2026-07-24 —
%     in review. ★ SPOTLIGHT MODULE. LOCKED-SPEC COMPLIANCE: mention-floor
%     phrasing throughout; NO rates (no denominator — instrument never
%     asked; ceiling restated in ¶1, M6 carries consequence); ACTIVATOR
%     stated SINGULAR (W26-R037/Q10 sole confirmation); non-activation =
%     evidenced pole (11); OPT-MANDATE 4 / OPT-CAP 13 / awareness-gap 12
%     (W25-R084 excluded per author ruling — genuine feature request);
%     gap = activation ∧ awareness; coordination-cost = the hinge, stated
%     as the survey's own arithmetic, theorizing DEFERRED (no module
%     labels in text). Disclosure: recount passes covered by M2.2 ¶2
%     footnote (option d) — no methods restatement here. Verbatims
%     verified char-exact vs open_text.csv 2026-07-24: W25-R038/Q10,
%     S25-R072/Q10, W25-R014/Q10, W25-R089/Q6, W26-R037/Q10, W26-R052/Q9,
%     W25-R060/Q10, S25-R040/Q10, W25-R071/Q9, W26-R022/Q10; S25-R043
%     instructor-session = paraphrase (abridged text only in 06 queue).
%     Platform facts = M2.1 DEP (restated in one clause, not re-described).
%     M3.6 COORDINATION NOTE: W26-R037 (the singular activator) is ALSO
%     F5's best within-respondent divergence pair AND an F1-PRESENCE row —
%     M3.6 should surface that continuity when it presents the pair.
%     Recount mechanics (115 candidates = 67 pre-refresh E-SHA ∪ 48-row
%     keyword sweep; E-SHA now 91 post-refresh; recount IMMUNE per Station
%     11) kept OUT of prose — floors are the adjudicated FINALs.
%     v2 RULINGS APPLIED 2026-07-24: ¶1 course-materials clause + the
%     AUTHOR-SUPPLIED first-person methods admission ("we did not think
%     to ask...falsely assumed...") VERBATIM per user 2026-07-24; floors
%     explained plainly (minimums; invisible non-mentioners); recount
%     sentence per user wording + four-kinds preview. ¶2 semicolon after
%     W25-R089/Q6 cite (user cited W25-R014 but the em-dash was at R089 —
%     applied there); "course design intended and which was heavily
%     encouraged" (user wording). ¶4 REBUILT: no cost/price/ledger/
%     arithmetic metaphors; "demand none of the solo uses make";
%     counts stated directly; "YouTube platform"; co-disclosedness/
%     world-disclosedness relation spelled out (students present-but-
%     alone, singly before monitors); plain closing sentence.

Of everything the virtual hospital offers, one affordance has no YouTube
counterpart at all: the world is built for watching together. Students
hear the classmates near them in the space, can share their screens
there, and every student belongs to a procedure-paired group from the
term's first week; the designed use of the environment is that the
group convenes inside the world and watches as a group, and that
multiplayer functionality is clearly stated and encouraged in the
course materials. What the survey can say about that design is limited
in a way the other findings are not. Unfortunately, we did not think to
ask the students whether or not their group convened in the world, so
no rate of use statistics are available. Given the group based nature
of this class, we falsely assumed students would organize meetings
within the VLE of their own volition, and that false assumption led the
course and survey designers to ignore crucial questions to be asked and
course designs that could have mitigated such. What the 873 responses
hold instead are incidental mentions---remarks students happened to
volunteer about watching together or about being alone---and counts
built from such mentions can only set minimums. That one student
described co-watching means at least one group convened; it cannot mean
only one did, because students who met and never mentioned it are
invisible to the survey. Each count in what follows is therefore a
floor---the minimum established by those who happened to speak---and
never a rate. After reviewing the relevant 115 candidate responses by
what each actually reports and sorting them, they fall into four
different kinds: reports that the group meeting never happened; a
report that it did; requests that the course structure the meeting; and
requests for a co-use capability that, in nearly every case, the
platform already possessed.

The best-attested kind is absence. Eleven students report directly that
the meeting did not happen for them. One logged in and found the
building empty: ``I found that when I did log in, I was usually alone
in the platform so it didn't really make me feel more connected with my
peers'' (W25-R038/Q10). One's group met, but elsewhere: ``The virtual
space to meet up for teamwork was very smart, but my group utilized
discord a lot'' (S25-R072/Q10). One knew the capability existed and
never used it: ``Meeting up with other people in the virtual space is
an interesting idea but because I didn't actually experience it all I
can say is it seems cool in theory'' (W25-R014/Q10). A fourth asked for
``more of a community since it just seemed isolating''
(W25-R089/Q6); a sentence that reads as a complaint about the design
until it is set beside the platform facts: the community infrastructure
existed, and what this student experienced was a groupmate-less
default, not an absence in the world. The isolation reported in these
responses is real. Its cause, in every case the responses let us see,
was that the assigned group never convened where the course design
intended and which was heavily encouraged.

Against these eleven stands one confirmation that the designed use
happened. In all 873 responses, exactly one student reports co-watching
as an experienced feature of the platform: ``It is more effective that
it does put you into a VR clinical environment and a cool interactive
space. The main pro to the platform is you can watch with others''
(W26-R037/Q10). No second response confirms it. The rest of the WE talk
is wishing, and the wishes split in two tellingly different ways. Four
students ask for the meeting to be structured---not for the capability,
but for someone to carry the organizing: ``Incorporating a required
group meetup within the VR simulation could incentivize more
interactive learning'' (W26-R052/Q9); a class sheet ``where times are
set by each group to enter the VR program'' (W25-R060/Q10); ``remote
teamwork VR space weekly meetings, to ensure accountability''
(S25-R040/Q10); an instructor logged in at the same time as the
students, explaining as they watch. Thirteen more ask for co-use
capability in wording that implies it does not exist---and for twelve
of the thirteen, the request is for something the platform already has.
``I think if there was a way to interact with other users who are
online at the same time that would add to it'' (W25-R071/Q9); ``perhaps
the ability to meet up with groups in a virtual reality setting''
(W26-R022/Q10)---the platform's existing capability, requested as an
innovation. The gap the survey shows is therefore double:
activation---groups that could have met and did not---and
awareness---students who never learned the meeting was possible at all.

Read together, the counts show why the meeting so rarely happened.
Every other use of the platform asks one student to act alone:
download, install, launch, walk, click play. Watching together is the
single use that requires two or more students to act at the same
time---agree on a time, log in together, find the same room---a demand
none of the solo uses make. The counts reflect that demand directly:
one student's responses confirm a group that met in the world; eleven
report groups that never did; and seventeen ask that the organizing be
done for them, by the course, the instructor, or a posted schedule. The
same pull of convenience that sent students to YouTube worked against
the meeting as well, and it worked hardest against the one affordance
the YouTube platform cannot reproduce: watching together, in one place,
at one time. In the dissertation's terms, what stood ready on this
route was co-disclosedness---the shared counterpart of
world-disclosedness. The students who described feeling present in the
virtual hospital were each alone in it: taken into the world, but
singly, one student before one monitor.
Co-disclosedness is that same being-taken-into-the-world had together
with others who are present in it at the same time, and it is the half
of the disclosure the platform's social architecture was built to make
possible. The survey states its standing directly: every student had it
available by design, and one student's responses confirm that it was
ever had.
% --- M3.6 From friction to justification (v1 DRAFTED 2026-07-24 — in
%     review. ANCHORS F5+F6. Trend-not-test stated FIRST (M2.2 DEP).
%     Sweep numbers canonical (02 §Emergent): 16 total, 5/2/9 by term =
%     1.3/0.7/4.2 per-100; sweep-methods transparency FOOTNOTED (artifact
%     code; 4 independent blind rediscoveries; circularity caveat).
%     Friction declines: crash 21→4→4, hw 8→4→1 (t3). A4-DIV 13
%     candidates / 12 unique (dup once). Residue = W26-R016/Q6 +
%     S25-R009/Q10 per F5. W26-R037 continuity surfaced (M3.5 activator =
%     divergence pair author); W26-R022 thread surfaced (skeptic proposes
%     the existing WE capability). Second-form gloss sourced from
%     construct-crosswalk §1 keystone vocabulary (occupied surface /
%     hollow depth; two-channel disagreement), NOT Part II v4, per
%     standing caution; NO German form-names (unverified); echoes-not-
%     detections ceiling stated. Verbatims verified char-exact
%     2026-07-24: W26-R033/Q9, W26-R056/Q9 (entire answer), W26-R022/Q10,
%     S25-R022/Q9, W25-R014/Q9 ("dont" typo verbatim), W26-R016/Q6,
%     S25-R009/Q10, W26-R037/Q9+Q10 ("then" typo verbatim). Hinge =
%     the justification question; theorizing deferred, no module labels.

The deployment's last finding lies in the difference between its
terms, and it must be read modestly: the three cohorts differ in size
and possibly in makeup, so what follows is a trend to be described,
not a test that was passed. The trend begins with the frictions
falling. Crash complaints numbered 21 in Winter 2025 and 4 in each
term after; hardware complaints fell from 8 to 4 to 1. By its third
term, the platform largely worked. What rose as the frictions fell was
a different kind of talk---talk that questioned the platform's value
rather than its function. A vocabulary-neutral sweep of all 873
responses finds 16 responses of this kind: five, two, and nine across
the three terms, which is 1.3\%, 0.7\%, and then 4.2\% of each term's
responses---the largest share by far in the final and smallest
cohort.\footnote{The counts come from a
keyword sweep applied uniformly across all three terms rather than from
the coded label for this theme, because that label was invented mid-run
by one batch's coders and its all-Winter-2026 distribution reflects
coder vocabulary, not the terms. The sweep found 16 responses (two
further keyword matches were inspected and discounted as false
positives). The theme itself is not an artifact: independent blind
coding passes rediscovered it in Winter 2025 under other names. One
caveat stands---the sweep's vocabulary is partly seeded from Winter
2026 phrasings, so the exact size of the rise is soft, though its
direction is robust.} The late voices are not reporting failure. ``I
think the youtube videos were enough to fully grasp the experience''
(W26-R033/Q9). ``Less dependency on the VR aspect altogether''
(W26-R056/Q9)---that is one student's entire answer to the improvement
question. Traditional learning, a third wrote, is ``the least
gimmicky'' way (W26-R022/Q10). The complaint has changed kind:
in the first term students told us the world did not work; by the
third, with the world working, they asked why it should exist.

A small register of responses sharpens that question by holding praise
and emptiness in one breath. Thirteen responses carry this pattern,
but two of them---submitted in different terms, under different
respondent codes---are word-for-word identical, and the analysis
counts that text once; the register therefore holds twelve distinct
responses that commend the platform and hollow it in the same breath. Read closely, most of them resolve
into ordinary complaints, because the hollowness has a named culprit.
``The actual idea of having the VR clinic is amazing, the videos and
images that come with it not so much'' (S25-R022/Q9)---the culprit is
the footage. ``The VR is a cool idea but the quality is so low I don't
gain from the atmosphere and I dont feel immersed'' (W25-R014/Q9)---the
culprit is the quality. Complaints of this shape are repairable in
principle: fix the named thing and the praise stands alone. Two
responses, though, name nothing. ``The virtual reality environment,
while it allows for more immersion, wasn't the highlight of the course,
and could be amended for something else'' (W26-R016/Q6). ``It is a cool
concept, but I feel like at the moment it is not doing much to help
regarding this class or the assignment'' (S25-R009/Q10). Nothing is
broken in these responses; something is missing, and the students
cannot say what. The sharpest case is a pair of answers by one student
---the same student whose response gave the survey its one confirmed
co-watching. Asked about educational use, this student praised the
platform in the survey's fullest presence language: it ``does put you
into a VR clinical environment,'' and ``you can watch with others''
(W26-R037/Q10). Asked what to improve, the same student wrote that the
platform ``is really laggy at times and is not significantly more
immersive then a YouTube video'' (W26-R037/Q9). One voice, both sides:
the world entered and shared, and the world flattened to a comparison
it barely survives.

What this register can testify to has a stated ceiling. In the
dissertation's account of boredom, which will be analyzed in the next
section, there is a form in which nothing determinate is
wrong---the occupation is intact, even pleasant, while the whole of it
quietly empties: an occupied surface over a hollow depth.\footnote{The
form is the second of the three Heidegger distinguishes in \textit{The
Fundamental Concepts of Metaphysics}: becoming bored \textit{with} an
occupation, as against being bored \textit{by} a determinate thing
(\textit{Fundamental Concepts} \S\S24--28). The full apparatus---the
three forms, their signatures, and what each requires of
evidence---is set out in the next section.} That form is detectable
only where two channels of evidence can disagree about the same
experience, as the laboratory's physiological record can disagree with
its behavioral one. A survey holds one channel: what a student
noticed, and chose to write, about their own experience. The
laboratory could set a brain measure against a gaze measure and let
the two contradict each other; a survey response has nothing set
against it that could contradict it. The two responses quoted above
that praise the platform while naming no culprit---W26-R016's and
S25-R009's---therefore show the same shape as this form of boredom:
commendation of an occupation that works, standing over an emptiness
the student cannot name. The shape matches; the survey cannot confirm
what lies beneath it. This section reads them as echoes of the form,
never as detections. What the survey can show directly is where the
arc of complaint ends. One further fact belongs to that ending: the
student who called the platform gimmicky went on, in the same
response, to propose its remedy---``perhaps the ability to meet up
with groups in a virtual reality setting'' (W26-R022/Q10)---the
platform's existing, unactivated capability. The survey's most
skeptical voice, answering a question about where virtual environments
could serve learning at all, proposes the one affordance this platform
offers that YouTube cannot: meeting inside the world, together. The
arc ends at a question no repair can answer: not whether the world
works, but what it is for---the justification the world itself must
supply.
